% TOP_PROBLEM: {"id": 6800001, "title": "Can you hear an orbifold singularity?", "category_id": 6, "category": {"id": 6, "name": "geometry", "display_name": "Geometry", "description": "Euclidean and non-Euclidean geometry, geometric structures.", "slug": "geometry", "order_index": 6, "created_at": "2026-07-31T15:26:25.671Z"}, "status": "partially_solved"}
% FIRST_POSTED: null
% ENTRY_KIND: statement_only
% Imported from: batch17.tex; UnsolvedMath ID 6800001
\documentclass[11pt]{article}
\usepackage{amsmath,amssymb,mathtools}
\usepackage{fontspec,unicode-math}
\setmainfont{Latin Modern Roman}
\setsansfont{Latin Modern Sans}
\setmathfont{Latin Modern Math}
\usepackage{xurl,hyperref}
\newcommand{\sref}[2]{\hyperlink{source:#1:#2}{[S#2]}}
\newcommand{\cataloguescope}[1]{\paragraph{Mathematical question.}#1}
\begin{document}
% UnsolvedMath ID 6800001; source catalogue code AMR-067-0001
\setcounter{section}{6800000}
\section{Can you hear an orbifold singularity?}\label{6800001}
{\small\sffamily UnsolvedMath ID 6800001 · Geometry}
\subsection{Definitions and mathematical statement}

\paragraph{Shared notation.}$\mathbb N=\{1,2,\ldots\}$, and $\mathbb Z,\mathbb Q,\mathbb R,\mathbb C$ denote the integers, rationals, reals and complex numbers. Any other conventions are specified locally.


A Riemannian orbifold is locally a quotient $\widetilde U/G_U$, where $\widetilde U\subset\mathbb R^n$ is open and connected, $G_U$ is a finite group acting effectively by diffeomorphisms, and a $G_U$-invariant Riemannian metric is compatible with the other charts. A point is singular if the stabilizer of a lift is nontrivial. An orbifold with no singular points is a smooth Riemannian manifold. Use the closed, compact and connected setting of the cited restatement. Two such orbifolds are Laplace-isospectral if the eigenvalues of the Laplace operator on smooth functions agree, including multiplicities.
\paragraph{Statement recorded in the supplied source.}
Do there exist Laplace-isospectral Riemannian orbifolds $O_1,O_2$ such that $O_1$ has a nonempty singular set and $O_2$ has no singular points? Equivalently, can the Laplace spectrum always detect the presence of an orbifold singularity? \sref{6800001}{2}
\subsection{Short English statement}
Can a compact orbifold with singular points and a smooth manifold have exactly the same function-Laplacian spectrum?
\subsection{Sources}
\begin{description}
\item[\hypertarget{source:6800001:1}{S1}] ULAM.AI and the UnsolvedMath Contributors, UnsolvedMath: A Curated Collection of Open Mathematics Problems, record 6800001 (AMR-067-0001). CC BY 4.0. \url{https://huggingface.co/datasets/ulamai/UnsolvedMath}.
\item[\hypertarget{source:6800001:2}{S2}] Sean Richardson and Elizabeth Stanhope, \emph{You Can Hear the Local Orientability of an Orbifold}, arXiv:1910.03224v1 (2019), Introduction, printed p. 2; Definition 2.1, pp. 2--3, and Laplace-spectrum convention, p. 4. \url{https://arxiv.org/abs/1910.03224v1}.
\end{description}
\label{end:6800001}
\end{document}
